\section*{Step 251: BN Harmonic Chain Extension to \(N=5000\)}

\paragraph{Arithmetic Gram formula.}
For the harmonic chain \(A_N=\{1/k:1\le k\le N\}\), steps 195 and 199 used
the arithmetic Gram identity
\[
G_{p,q}=
\sum_{n\ge 1}
\frac{((n\bmod p)/p)((n\bmod q)/q)}{n(n+1)}.
\]
Step 251 extends the same pipeline to \(N=3000,4000,5000\).  The truncated
series used \(X_{\max}=300000\), giving per-entry tail bound
\[
X_{\max}^{-1}=3.33333\cdot 10^{-6}.
\]
Scalar setup used mpmath at 80 dps; the active matrices were evaluated by
symmetric eigen-pseudoinverse with cutoff \(10^{-12}\lambda_{\max}\).

\paragraph{Results.}
\[
\begin{array}{c|c|c|c}
N & \delta_A^2 & \delta_A^2\log N & \text{condition proxy}\\
\hline
3000 & 0.00558973365837 & 0.0447534622742 & 4.48\cdot10^7\\
4000 & 0.00538148417111 & 0.0446342968526 & 8.69\cdot10^7\\
5000 & 0.00525837798558 & 0.0447866211767 & 1.42\cdot10^8
\end{array}
\]
The \(N=2000\) calibration under the same \(X_{\max}=300000\) truncation gives
\[
\delta_A^2(2000)=0.00609996484863,\qquad
\delta_A^2(2000)\log(2000)=0.0463652378211.
\]

\paragraph{Fit.}
The tail mean over \(N\ge 200\) through \(5000\) is
\[
\operatorname{mean}(\delta_A^2\log N)=0.045311325284006.
\]
The recent tail mean over \(N\ge1000\) is
\[
0.0448554780938519.
\]
A fit of the form \(C/(\log N)^\alpha\) on \(N\ge1000\) gives
\[
C=0.0488745892509167,\qquad \alpha=1.041904016.
\]

\paragraph{Verdict.}
The data are consistent with the slow logarithmic Baez--Duarte/Nyman--Beurling
regime through \(N=5000\).  This is a finite numerical diagnostic, not a proof
of RH or of convergence of the infinite criterion.
